\documentclass[12pt,a4paper]{article}

\usepackage{fontspec}
\usepackage{xeCJK}
\usepackage[a4paper,top=3cm,bottom=2.54cm,left=3.175cm,right=3.175cm,
            footskip=1.15cm]{geometry}
\usepackage{fancyhdr}
\usepackage{caption}
\usepackage{amsmath,amssymb}
\usepackage{array,tabularx,longtable}
\usepackage{enumitem}
\usepackage{tikz}
\usepackage{float}
\usepackage{xcolor}
\usetikzlibrary{arrows.meta,positioning,shapes.geometric,fit}

% 中文统一为宋体，西文和数字统一为 Times New Roman。
\setmainfont{Times New Roman}
\setmonofont{Times New Roman}
\setCJKmainfont[AutoFakeBold=2.5]{SimSun}
\setCJKsansfont[AutoFakeBold=2.5]{SimSun}
\setCJKmonofont{SimSun}

\newcommand{\xiaosi}{\fontsize{12bp}{22bp}\selectfont}
\newcommand{\wuhao}{\fontsize{10.5bp}{15bp}\selectfont}
\newcommand{\placeholder}{\textcolor{gray}{待实测}}

\setlength{\parindent}{2em}
\setlength{\parskip}{0pt}
\setlength{\tabcolsep}{4pt}
\renewcommand{\arraystretch}{1.35}
\setlist{nosep}

\pagestyle{fancy}
\fancyhf{}
\renewcommand{\headrulewidth}{0pt}
\fancyfoot[R]{\wuhao\thepage}

\captionsetup{
  font={small},
  labelfont=normalfont,
  labelsep=space,
  justification=centering,
  singlelinecheck=true,
  skip=3pt
}
\renewcommand{\figurename}{图}
\renewcommand{\tablename}{表}

\newcommand{\mainheading}[1]{%
  \par\vspace{7.8pt}\noindent{\fontsize{16bp}{22bp}\selectfont\bfseries #1}\par
  \vspace{4pt}}
\newcommand{\subheading}[1]{%
  \par\vspace{2pt}\noindent{\fontsize{14bp}{22bp}\selectfont\bfseries #1}\par}
\newcommand{\thirdheading}[1]{%
  \par\noindent{\xiaosi\bfseries #1}\par}

\tikzset{
  block/.style={draw,align=center,minimum height=8mm,minimum width=20mm},
  decision/.style={draw,diamond,aspect=2,align=center,inner sep=1pt},
  line/.style={-{Latex[length=2mm]},thin}
}

\begin{document}
\xiaosi

% 参考报告中摘要单独成页，不显示页码。
\thispagestyle{empty}
\begin{center}
  {\fontsize{16bp}{22bp}\selectfont\bfseries 摘要}
\end{center}

本工程由电流互感取能与采样模块、SS16 桥式整流及稳压模块、可编程增益测量模块、
MSPM0G3507 控制模块、本地显示与蓝牙通信模块、Android 无线读表器组成。电流互感器
同时获取工作能量和测量信号，控制器完成自动量程、真有效值计算及外设电源管理，
OLED 显示本地电流，蓝牙和 Android 读表器完成远程读取、报警、记录与趋势显示。

\noindent\textbf{关键词：}无源电流表；MSPM0G3507；真有效值；自动量程；低功耗；BLE

\clearpage
\setcounter{page}{1}

\mainheading{一、\quad 系统方案论证与比较}

\subheading{1．取能与传感方案的论证与选择}

方案一：使用两组独立次级线圈分别驱动取电电路和传感电路。该方案中两条通道互相
隔离，调试较为直观，但两组绕组同时占用磁芯窗口，绕制工作量较大，匝数和耦合差异
还会引入一致性误差。

方案二：使用一组带抽头的次级线圈，取能与传感共用整流链路，再从统一节点分配供电
和测量信号。由于两条通道都需要对交流量进行整流，共用方案可以减少重复的线圈、
整流器件和导通损耗，并通过抽头调整供电电压与测量动态范围。

综上所述，本系统采用方案二。初版绕组为 85 匝，后续改为总计 270 匝、每 50 匝
预留一个抽头的结构。经不同抽头组合测试，最终以 170 匝有效绕组为取电电路供电，
其余抽头保留用于测量与后续参数调整。

\subheading{2．整流及稳压电源方案的论证与选择}

普通硅二极管桥式整流结构简单，但在低输入电压时，两个串联导通二极管的正向压降
占比较大，会明显降低最低电流下的可用能量。肖特基二极管具有较低的正向导通压降，
更适合本题的低压取能场景。因此本系统采用四只 SS16 构成桥式整流。

整流后的稳压可采用 DC-DC 或 LDO。DC-DC 在输入输出压差较大时具有较高的能量利用率，
但开关纹波、启动电压、外围器件和控制器静态功耗均需重点评估；LDO 外围简单、噪声
较低，有利于模拟测量，但效率近似为输出电压与输入电压之比，压差较大时损耗较高。
两种方案比较如表 \ref{tab:power_compare} 所示。

\begin{table}[H]
\centering
\caption{稳压电源方案比较}
\label{tab:power_compare}
\wuhao
\begin{tabularx}{\textwidth}{|>{\centering\arraybackslash}p{23mm}|
>{\centering\arraybackslash}X|>{\centering\arraybackslash}X|}
\hline
比较项目 & DC-DC 稳压 & LDO 稳压\\ \hline
能量利用率 & 压差较大时通常更高 & 约为 $V_{\mathrm{o}}/V_{\mathrm{i}}$\\ \hline
启动条件 & 取决于芯片启动电压和外围参数 & 输入需高于输出及压差要求\\ \hline
纹波与噪声 & 存在开关纹波，需滤波和布局优化 & 输出纹波较小，利于模拟测量\\ \hline
器件与调试 & 外围器件较多，参数设计较复杂 & 外围少，调试简单\\ \hline
最终选型 & \multicolumn{2}{c|}{\placeholder（根据最低启动电流、效率和纹波测试填写）}\\ \hline
\end{tabularx}
\end{table}

\subheading{3．电流有效值测量及量程方案的论证与选择}

方案一：使用峰值检波电路测量交流峰值，再按理想正弦关系换算有效值。该方案运算量
小，但二极管压降和保持电容误差在 0.1\,A 附近较明显，并且被测波形一旦偏离正弦，
换算结果便会产生附加误差。

方案二：直接对波形采样并计算真有效值。该方案可在数字域内消除直流偏置，并减小
波形失真和采样记录非整数周期的影响。MSPM0G3507 具备 ADC、定时器和 DMA 资源，
能够完成 4\,kHz 连续采样，故本设计选择方案二，并使用 Hann 权重计算真有效值。

固定高增益虽能提高小信号分辨率，但大电流时容易饱和；固定低增益则使小电流所占
ADC 码宽过小。因此模拟前端采用 1、2、4、8、16、32 倍六档自动量程，优先选择
不饱和条件下的最高增益，并在相邻档之间设置回差以避免门限附近反复切换。

\mainheading{二、\quad 理论分析与计算}

\subheading{1．电流互感器与绕组参数分析}

电流互感器原边串入被测交流线路，原边匝数满足 $N_1\leq5$，副边匝数为 $N_2$。
忽略励磁电流与损耗时，原、副边安匝近似平衡，即
\begin{equation}
  N_1 I_1=N_2 I_2 ,
\end{equation}
因此副边电流为
\begin{equation}
  I_2=\frac{N_1}{N_2}I_1 .
\end{equation}
增加副边匝数能够提高低电流时的感应电压，但同时会增加绕组电阻和分布参数。为兼顾
取能和调试，本设计由 85 匝初版改为 270 匝带抽头结构。试验时分别比较不同有效匝数
下的整流电压、带载稳定性和启动能力，最终选取 170 匝有效绕组供电。

\subheading{2．桥式整流与储能分析}

桥式整流每半周均有两只二极管串联导通。若互感器次级瞬时电压为 $u_2(t)$，单只
SS16 的正向压降为 $V_{\mathrm{F}}$，则整流后的近似瞬时电压为
\begin{equation}
  u_{\mathrm{r}}(t)=\max\left(|u_2(t)|-2V_{\mathrm{F}},\,0\right).
\end{equation}
由式（3）可知，输入电压越低，二极管压降对可用能量的影响越明显，故选用肖特基桥。
储能电容需要同时满足 MCU 完成一次探测、正式采样、计算及外设开启过程中的能量需求。
最低启动电流应在电容完全放电的条件下实测，不能仅由空载整流电压判断。

\subheading{3．可编程增益与偏置计算}

对增益为 $G$ 的同相偏置放大电路，其输出可表示为
\begin{equation}
  V_{\mathrm{o}}=G V_{\mathrm{i}}+(1-G)V_{\mathrm{DAC}} .
\end{equation}
若输入直流中心为 $V_{\mathrm{i,dc}}$，希望输出中心位于 ADC 中点 $V_{\mathrm{c}}$，
则 $G>1$ 时的 DAC 电压为
\begin{equation}
  V_{\mathrm{DAC}}=\frac{G V_{\mathrm{i,dc}}-V_{\mathrm{c}}}{G-1}.
\end{equation}
固件以 12 位 ADC 的 2047 码为目标中心，探测阶段采集 160 点，时间为 40\,ms；
正式阶段采集 800 点，时间为 200\,ms。目标峰峰值利用率设为满量程的 65\%，回差
下、上限分别取 25\% 和 75\%。

\subheading{4．真有效值及标定算法}

定时器 TIMG6 以 4\,kHz 产生 ADC1 转换事件，DMA 将数据搬运至内存。对 $N$ 点
采样序列 $x[n]$ 使用 Hann 权重 $w[n]$，其加权均值为
\begin{equation}
  \bar{x}_{w}=
  \frac{\sum_{n=0}^{N-1}w[n]x[n]}{\sum_{n=0}^{N-1}w[n]} .
\end{equation}
去除直流分量后的有效值码宽为
\begin{equation}
 X_{\mathrm{rms}}=
 \sqrt{\frac{\sum_{n=0}^{N-1}w[n]\left(x[n]-\bar{x}_{w}\right)^2}
 {\sum_{n=0}^{N-1}w[n]}} .
\end{equation}
Hann 权重可降低采样记录没有覆盖整数工频周期时的边界影响。各量程分别保存标定点
$(x_k,I_k)$。当 $x_k\leq X_{\mathrm{rms}}\leq x_{k+1}$ 时，采用
\begin{equation}
 I=I_k+\frac{X_{\mathrm{rms}}-x_k}{x_{k+1}-x_k}(I_{k+1}-I_k)
\end{equation}
进行分段线性插值。当前固件已经保留六档标定表接口；正式验收前必须写入实测标定点，
名义比例系数只可用于调试，不能作为 0.5\% 精度的证明。

\mainheading{三、\quad 电路与程序设计}

\subheading{1．系统框图}

系统由电流互感器、SS16 桥式整流、储能与稳压、可编程模拟前端、MSPM0G3507、
OLED、受控供电的 HC-42 以及 Android 读表器组成。电流表 A、B 结构相同且供能与
测量链路相互独立，拆除任意一表不会影响另一表。硬件系统框图如图
\ref{fig:system_v2} 所示。

\begin{figure}[H]
\centering
\begin{tikzpicture}[node distance=6mm and 6mm,scale=0.74,transform shape]
  \node[block] (line) {36\,V 交流线路\\0.1～2.0\,A};
  \node[block,right=of line] (ct) {CT\\270 匝带抽头};
  \node[block,right=of ct] (bridge) {SS16\\桥式整流};
  \node[block,right=of bridge] (power) {储能/稳压\\模拟前端};
  \node[block,below=of power] (mcu) {MSPM0G3507\\RMS/量程/地址};
  \node[block,left=of mcu] (oled) {S2 按键\\OLED};
  \node[block,right=of mcu] (ble) {PB17 电源控制\\HC-42};
  \node[block,right=of ble] (app) {Android 读表器\\报警/存储/趋势};
  \draw[line] (line)--(ct);
  \draw[line] (ct)--(bridge);
  \draw[line] (bridge)--(power);
  \draw[line] (power)--(mcu);
  \draw[line] (mcu)--(oled);
  \draw[line] (mcu)--(ble);
  \draw[line] (ble)--(app);
  \node[draw,dashed,fit=(ct)(bridge)(power)(mcu)(oled)(ble),
        label={[font=\wuhao]above:电流表 A（电流表 B 结构相同且独立）}] {};
\end{tikzpicture}
\caption{硬件系统框图}
\label{fig:system_v2}
\end{figure}

\subheading{2．取能、整流与稳压电路设计}

取能绕组采用 270 匝总绕组并设置抽头，最终选用其中 170 匝有效绕组为整流桥供电。
四只 SS16 组成全桥，使正、负半周均向储能电容充电。储能节点再经稳压电路为 MCU
和模拟前端供电。DC-DC 与 LDO 的最终选型应依据最低启动电流、带载效率和输出纹波
测试确定，原理图中应同时标出储能电容放电通路，便于每次启动试验前恢复相同初始状态。

\subheading{3．模拟前端与控制器接口设计}

模拟前端提供直通 X1 及 OPA 可编程增益 X2、X4、X8、X16、X32 六档。DAC 根据
探测阶段得到的直流中心和当前增益调整偏置，避免高增益时输出中心偏离 ADC 中点。
MSPM0G3507 的 ADC、OPA、OLED、按键、蓝牙电源、UART 和地址接口分配见附录表
\ref{tab:pins_v2}。

\subheading{4．电流测量程序设计}

程序上电后先把 PB17 拉低并保持 OLED 未初始化。每次测量先采集 160 点探测记录，
根据峰峰值和直流中心选择增益及 DAC，再采集 800 点正式记录，通过 Hann 加权计算
真有效值并查找对应量程的标定表。程序流程如图 \ref{fig:flow_v2} 所示。

\begin{figure}[H]
\centering
\begin{tikzpicture}[node distance=5mm and 20mm,scale=0.84,transform shape]
  \node[block] (boot) {上电\\PB17=0，OLED 不初始化};
  \node[block,below=of boot] (probe) {160 点探测\\选择增益和 DAC};
  \node[block,below=of probe] (sample) {800 点 DMA 采样\\计算 Hann 真有效值};
  \node[decision,below=of sample] (first) {首次测量\\完成？};
  \node[block,right=of first] (bt) {PB17=1\\等待 350\,ms\\启动 UART};
  \node[block,below=of first] (send) {形成数据帧\\周期发送};
  \node[decision,below=of send] (key) {S2 按下？};
  \node[block,right=of key] (oled) {切换 OLED 状态\\仅显示电流有效值};
  \node[block,below=of key] (next) {进入下一测量周期};
  \draw[line] (boot)--(probe);
  \draw[line] (probe)--(sample);
  \draw[line] (sample)--(first);
  \draw[line] (first)--node[left]{否}(send);
  \draw[line] (first.east)--node[above]{是}(bt.west);
  \draw[line] (bt.south)|-(send.east);
  \draw[line] (send)--(key);
  \draw[line] (key)--node[left]{否}(next);
  \draw[line] (key.east)--node[above]{是}(oled.west);
  \draw[line] (oled.south)|-(next.east);
\end{tikzpicture}
\caption{电流测量与外设控制流程图}
\label{fig:flow_v2}
\end{figure}

\subheading{5．OLED 与蓝牙电源控制设计}

OLED 上电时不发送初始化命令。S2 采用状态切换逻辑：由关闭状态单击一次后才初始化
并持续显示电流有效值，再次单击则关闭显示，用户无需一直按住按键。OLED 对比度命令
设置为 \texttt{0x81,0x0F}。

蓝牙模块不是仅延后 UART，而是由 PB17 控制外部 3.3\,V 负载开关。PB17 低电平时
负载开关关闭，高电平时导通。首次完整测量完成后 PB17 才置高，等待 350\,ms 供电
稳定，再初始化 UART。蓝牙断电期间 UART 引脚保持未初始化，避免串口反向灌电。
仅计算固件固定时序，蓝牙最早约在 $40+200+350=590$\,ms 后具备通信条件，实际启动
时间还包括互感取能和储能电压建立时间。

\subheading{6．无线通信与 Android 读表器设计}

电流表按 ASCII 行发送：
\begin{center}
{\wuhao\texttt{METER\_TEST,ADDR=01,CURRENT\_MA=1234,STATUS=NORMAL\textbackslash r\textbackslash n}}
\end{center}
地址范围为 00～15，电流单位为 mA。HC-42 使用 FFE0 服务和 FFE1 通知特征。
Android 应用启动后扫描 8\,s，优先连接名称符合 HC-42、HC42、METER、AMMETER 或
BYJX\_ 前缀的设备；基本部分提供人工触发一次读取，发挥部分默认每 120\,s 自动搜索。
SQLite 数据库保存 1000 条读数，历史列表和趋势显示均超过题目要求的 30 条。当电流
低于 0.2\,A、超过 2.0\,A 或设备离线时分别产生低限、超限和离线提示。

\mainheading{四、\quad 测试方案与测试结果}

\subheading{1．系统测试方案}

\thirdheading{（1）测试仪器}

测试仪器计划如表 \ref{tab:instruments_v2} 所示，具体型号以最终实验室设备为准。

\begin{table}[H]
\centering
\caption{测试仪器列表}
\label{tab:instruments_v2}
\wuhao
\begin{tabularx}{\textwidth}{|>{\centering\arraybackslash}p{10mm}|
>{\centering\arraybackslash}p{32mm}|>{\centering\arraybackslash}p{27mm}|
>{\centering\arraybackslash}p{10mm}|X|}
\hline
序号 & 仪器类别 & 仪器型号 & 数量 & 性能参数及用途\\ \hline
1 & 36\,V 交流隔离电源 & \placeholder & 1 & 容量不小于 60\,VA\\ \hline
2 & 可调负载/电流源 & \placeholder & 1 & 产生 0.1～2.0\,A 稳定交流电流\\ \hline
3 & 标准交流电流表 & \placeholder & 1 & 准确度优于被测表至少 3 倍\\ \hline
4 & 数字示波器 & \placeholder & 1 & 观察储能、采样及启动时序\\ \hline
5 & Android 手机 & \placeholder & 1 & 5\,m 通信、报警和存储测试\\ \hline
\end{tabularx}
\end{table}

\thirdheading{（2）测试方法}

1、进行电流测量精度测试。

（1）保持线路电压为 36\,V，依次设置 0.10、0.20、0.50、1.00、1.50 和
2.00\,A，每点稳定后记录标准表与被测表读数。

（2）每个电流点重复测试不少于 5 次，在量程重叠区增加测试点，并分别完成 A、B
两表测试。

（3）拆除电流表 B 后重复测试 A，比较误差变化，验证两表互不影响。

2、进行取能与启动性能测试。

（1）储能电容充分放电后，从较小电流逐步增加，记录能够稳定完成测量和发送的最低
电流。

（2）在代表电流点捕获储能电压建立、首次有效值生成和首帧到达时刻。

（3）比较不同抽头组合的整流电压和带载稳定性，记录最终供电匝数。

3、进行本地显示和无线功能测试。

（1）确认上电阶段 OLED 无驱动；单击 S2 后持续显示，再次单击关闭。

（2）手机与电流表间距设为 5\,m，分别检查人工读取和 120\,s 周期搜索。

（3）设置低限、正常、超限并中断一只表通信，检查报警；写入不少于 30 条记录，
重启应用后检查历史和趋势数据。

\subheading{2．测试结果与分析}

\thirdheading{（1）供电绕组测试}

初版次级绕组为 85 匝，后续绕制为总计 270 匝并设置抽头。通过不同抽头组合对比
整流输出和带载稳定性后，最终选取 170 匝有效绕组供电。该结果已经用于当前取能方案。

\begin{table}[H]
\centering
\caption{供电绕组方案演进}
\wuhao
\begin{tabularx}{\textwidth}{|>{\centering\arraybackslash}p{28mm}|
>{\centering\arraybackslash}p{35mm}|X|}
\hline
阶段 & 绕组结构 & 结果\\ \hline
初版 & 85 匝次级绕组 & 感应电压和调节余量有限\\ \hline
改进版 & 270 匝总绕组，每 50 匝预留抽头 & 可组合不同有效匝数进行测试\\ \hline
最终方案 & 170 匝有效绕组供电 & 经抽头对比测试确定\\ \hline
\end{tabularx}
\end{table}

\thirdheading{（2）电流测量精度测试}

被测显示值 $I_{\mathrm{x}}$ 相对标准值 $I_{\mathrm{ref}}$ 的相对误差为
\begin{equation}
  \delta=
  \frac{I_{\mathrm{x}}-I_{\mathrm{ref}}}{I_{\mathrm{ref}}}\times100\% .
\end{equation}
题目要求 0.1～2.0\,A 全量程内满足 $|\delta|\leq0.5\%$。当前尚未完成标准表对比和
六档实物标定，故表 \ref{tab:accuracy_v2} 保留为测试记录模板，不填入推测数据。

\begin{table}[H]
\centering
\caption{电流测量精度记录表}
\label{tab:accuracy_v2}
\wuhao
\begin{tabular}{|c|c|c|c|c|c|c|}
\hline
标准值/A & 0.10 & 0.20 & 0.50 & 1.00 & 1.50 & 2.00\\ \hline
A 表示值/A & \placeholder & \placeholder & \placeholder & \placeholder & \placeholder & \placeholder\\ \hline
A 表误差/\% & \placeholder & \placeholder & \placeholder & \placeholder & \placeholder & \placeholder\\ \hline
B 表示值/A & \placeholder & \placeholder & \placeholder & \placeholder & \placeholder & \placeholder\\ \hline
B 表误差/\% & \placeholder & \placeholder & \placeholder & \placeholder & \placeholder & \placeholder\\ \hline
\end{tabular}
\end{table}

\thirdheading{（3）阶段性功能检查}

固件格式、编译和静态检查已经通过，OLED 按键切换、PB17 蓝牙受控供电和 Android
数据链路已经实现；精度、最低启动电流、5\,m 距离及双表报警仍需实物联调。完整
检查表见附录表 \ref{tab:check_v2}。

现阶段可以确认程序控制逻辑和通信链路已经闭环，但名义比例系数不能证明 0.5\%
测量精度。预计主要误差来源包括互感器变比和相位误差、采样电阻温漂、运放增益误差、
ADC 参考误差以及档位交界处不连续。后续应在 0.1～0.3\,A 区间加密标定点，并对
相邻量程重叠区进行一致性修正。

\mainheading{五、\quad 总结}

本系统完成了``无源''交流电流表及无线读表器的总体方案、测量程序和手机端软件设计。
取能与传感共用带抽头的次级绕组和 SS16 肖特基整流桥，抽头对比测试最终确定以
170 匝有效绕组供电；六档自动量程、DAC 偏置、定时采样、DMA 和加窗真有效值算法
提高了 ADC 量程利用率。OLED 按需初始化，蓝牙采用 PB17 控制外部电源，减小了上电
阶段的瞬态负担。Android 读表器能够完成地址识别、人工读取、周期搜索、异常报警、
历史记录和趋势显示。当前不足是精度、最低启动电流、启动时间、5\,m 距离及双表
联调尚缺实测数据，后续应按测试表完成全量程标定并以实测结果替换占位项。

硬件方面还需在相同互感器、储能电容和负载条件下比较 DC-DC 与 LDO 的最低启动电流、
稳定效率和输出纹波，确定正式稳压方案；同时完善储能节点放电通路，使每次冷启动测试
具有一致的初始条件。170 匝供电结论应进一步在 0.1～2.0\,A 全范围内检查稳压余量，
避免仅以某一工作点的整流电压作为选型依据。

标定方面应分别为 X1～X32 六档采集多组标准点，在 0.1～0.3\,A 区间及相邻量程
重叠区加密测试，并比较温升前后的误差。系统联调时应同时接入 A、B 两表，逐项验证
地址设置、人工读取、120\,s 周期搜索、低限/超限/离线报警、30 条以上记录保存及
5\,m 通信距离。完成上述工作后，系统才能形成用于正式评审的完整测试结论。

\clearpage
\mainheading{附录}

\subheading{1．题目指标落实对照}

\begin{longtable}{|p{54mm}|p{82mm}|}
\caption{题目要求与设计落实情况}\\
\hline
题目要求 & 设计落实\\ \hline
\endfirsthead
\hline
题目要求 & 设计落实\\ \hline
\endhead
仅由电磁耦合供电，不能外接电源 & CT 取能，经 SS16 桥式整流、储能和稳压后供电；
最终需实物验收\\ \hline
原边匝数不超过 5 匝 & 结构设计约束 $N_1\leq5$，副边总计 270 匝并设置抽头\\ \hline
0.1～2.0\,A，误差不超过 0.5\% & 六档自动量程与分段标定；精度待标准表实测\\ \hline
A、B 两表独立，移除 B 不影响 A & 两表使用独立互感器、储能和测量链路\\ \hline
4 位地址和串行接口 & PB0、PB6～PB8 设置低有效地址；UART2 连接 HC-42\\ \hline
具有放电功能 & 测试前通过储能回路放电；硬件放电方式需在原理图和实物标注\\ \hline
无线距离不小于 5\,m & HC-42 与 Android 手机通信；距离待现场验证\\ \hline
人工启动 1 次读取 & App 提供人工触发下一次读取\\ \hline
每 2\,min 自动搜索 & App 默认 120\,s 自动扫描并保存周期设置\\ \hline
低限、超限和离线报警 & 阈值为 0.2\,A、2.0\,A，超时判定离线\\ \hline
不少于 30 条记录、掉电保存和趋势 & SQLite 保存 1000 条，历史和趋势显示均超过 30 条\\ \hline
外形不超过 $12\times7\times5$\,cm & 作为 PCB、外壳和最终验收约束\\ \hline
\end{longtable}

\subheading{2．控制器接口与阶段性检查}

\begin{table}[H]
\centering
\caption{MSPM0G3507 主要接口分配}
\label{tab:pins_v2}
\wuhao
\begin{tabularx}{\textwidth}{|>{\centering\arraybackslash}p{25mm}|
>{\centering\arraybackslash}p{30mm}|X|}
\hline
引脚 & 功能 & 说明\\ \hline
PA18 & ADC1 输入 & 电流采样原始通道\\ \hline
PA16 & 模拟输入 & OPA 可编程增益输出通道\\ \hline
PB2、PB3 & 软件 I\textsuperscript{2}C & OLED 的 SCL、SDA\\ \hline
PB21 & S2 按键 & 下降沿切换显示状态，松手后防抖\\ \hline
PB17 & 蓝牙电源使能 & 高电平导通外部负载开关，低电平关闭\\ \hline
PB15、PB16 & UART2 & HC-42 串口，9600\,bit/s\\ \hline
PB0、PB6～PB8 & 地址输入 & 4 位低有效地址，范围 00～15\\ \hline
\end{tabularx}
\end{table}

\begin{table}[H]
\centering
\caption{系统功能与指标检查表}
\label{tab:check_v2}
\wuhao
\begin{tabularx}{\textwidth}{|p{40mm}|>{\centering\arraybackslash}p{22mm}|
>{\raggedright\arraybackslash}X|}
\hline
检查项目 & 当前状态 & 说明\\ \hline
固件格式、编译和静态检查 & 已完成 & \texttt{cargo fmt}、\texttt{cargo check} 和
\texttt{cargo clippy -D warnings} 均通过\\ \hline
OLED 上电关闭、按键切换 & 已实现 & S2 单击切换，无需持续按住\\ \hline
蓝牙受控供电 & 已实现 & PB17 低关闭、高导通，首次测量后开启\\ \hline
供电绕组抽头对比 & 已完成 & 最终选取 170 匝有效绕组供电\\ \hline
0.1～2.0\,A、误差不超过 0.5\% & \placeholder & 需完成各档实物标定\\ \hline
最低启动电流和启动时间 & \placeholder & 需在完全放电后测量\\ \hline
5\,m 无线距离、双表和报警 & \placeholder & 需连接两只实物完成联调\\ \hline
外形尺寸限制 & \placeholder & 装配外壳后测量\\ \hline
\end{tabularx}
\end{table}

\subheading{3．通信协议}

\noindent 帧格式：
\begin{center}
\texttt{METER\_TEST,ADDR=dd,CURRENT\_MA=iiii,STATUS=ssss\textbackslash r\textbackslash n}
\end{center}

\noindent 字段说明：

（1）\texttt{ADDR}：十进制两位电表地址，00～15。

（2）\texttt{CURRENT\_MA}：电流有效值，单位为 mA。

（3）\texttt{STATUS}：\texttt{NORMAL}、\texttt{LOW}、\texttt{HIGH} 或
\texttt{OFFLINE}。

\clearpage
\subheading{4．分档标定记录}

\begin{table}[H]
\centering
\caption{分档标定原始记录表}
\wuhao
\begin{tabular}{|c|c|c|c|c|c|}
\hline
量程 & 标准电流/A & RMS 码宽 & 显示电流/A & 误差/\% & 温度/$^\circ$C\\ \hline
X32 & 0.10 & & & &\\ \hline
X32/X16 & 0.20 & & & &\\ \hline
X16/X8 & 0.30 & & & &\\ \hline
X8/X4 & 0.50 & & & &\\ \hline
X4/X2 & 1.00 & & & &\\ \hline
X2/X1 & 1.50 & & & &\\ \hline
X1 & 2.00 & & & &\\ \hline
\end{tabular}
\end{table}

\subheading{5．最终检查项目}

（1）确认中文正文、标题和图表题全部使用宋体，英文和数字使用 Times New Roman。

（2）确认摘要页无页码，正文从 1 开始且页码位于右下角。

（3）确认 A、B 两表独立，地址设置正确，PB17 蓝牙电源逻辑为低关、高开。

（4）确认正式提交前已用实测标定表替换调试比例系数，并填写全部待测数据。

\end{document}
